\documentclass[../../../include/open-logic-section]{subfiles}

\begin{document}

% Part: turing-machines 
% Chapter: undecidability 
% Section: verification

\olfileid{tur}{und}{ver}
\olsection{验证表示}

\begin{explain}
为了验证我们的表示是可行的，我们必须证明两件事。首先，我们必须证明，如果 $M$ 在输入~$w$ 上停机，那么
$!T(M, w) \lif !E(M, w)$ 是有效的。其次，我们必须证明其逆命题，
即，如果 $!T(M, w) \lif !E(M, w)$ 是有效的，那么 $M$ 在输入~$w$ 上运行时最终确实会停机。

证明这两点的策略非常不同。对于第一个结果，我们必须证明一个一阶逻辑语句（即
$!T(M, w) \lif !E(M, w)$）是有效的。最简单的办法是给出一个推导。然而，我们的证明需要对所有
$M$ 和 $w$ 都成立，因此我们并非只需为单个语句给出推导，而是要为无穷多个语句给出推导。
所以，我们能采取的最佳策略是用归纳法证明：无论 $M$ 和~$w$ 具体是什么样，也无论 $M$ 在输入~$w$ 上停机需要多少步，总存在 $!T(M, w) \lif !E(M, w)$ 的推导。

很自然地，我们的归纳将基于 $M$ 抵达停机配置之前所用的步数进行。在我们的归纳证明中，将确立：对于 $M$ 在输入~$w$ 上运行的每一步~$n$，
$!T(M, w) \Entails !C(M, w, n)$，其中 $!C(M, w, n)$ 正确地
描述了 $M$ 在~$w$ 上运行~$n$ 步后的配置。那么，如果
$M$ 在输入~$w$ 上，比如说，于 $n$ 步后停机，$!C(M, w, n)$ 将描述一个停机配置。我们还将证明，只要 $!C(M, w, n)$ 描述一个停机配置，
$!C(M, w, n)
\Entails !E(M, w)$。因此，如果 $M$ 在输入~$w$ 上停机，则存在某个~$n$，使得 $M$
在 $n$~步后处于一个停机配置。于是，
$!T(M, w)
\Entails !C(M, w, n)$，其中 $!C(M, w, n)$ 描述了一个停机配置，且由于在该情况下
$!C(M, w, n) \Entails !E(M, w)$，
我们得到 $T(M, w) \Entails !E(M, w)$，即
$\Entails !T(M, w)
\lif !E(M, w)$。

证明其逆命题的策略则迥然不同。这里我们假设
$\Entails !T(M, w) \lif !E(M, w)$，并且必须证明 $M$ 在输入~$w$ 上停机。
由该假设我们得到 $!T(M, w) \Entails !E(M,
w)$，即 $!E(M, w)$ 在所有使得 $!T(M,
w)$ 为真的结构中为真。因此，我们将描述一个结构~$\Struct{M}$，在其中
$!T(M, w)$ 为真：其论域将是 $\Nat$，且所有 $\Obj Q_q$ 和 $\Obj S_\sigma$ 的解释将由 $M$ 在输入~$w$ 上运行期间的配置给出。举例来说，
$\Sat{M}{\Obj Q_q(\num{m}, \num{n})}$ 当且仅当 $M$ 在输入~$w$
上运行 $n$ 步后，正处于状态~$q$ 且扫描着方格~$m$。现在，由于根据假设
$!T(M, w) \Entails !E(M, w)$，并且根据构造有
$\Sat{M}{!T(M,
  w)}$，因此
$\Sat{M}{!E(M, w)}$。 但是
$\Sat{M}{!E(M, w)}$
当且仅当存在某个 $n \in \Domain{M} = \Nat$，使得 $M$ 在输入~$w$ 上运行 $n$ 步后处于一个停机配置。
\end{explain}

\begin{defn}
令 $!C(M, w, n)$ 为如下语句：
\[ 
\Obj Q_q(\num{m}, \num{n}) \land \Obj S_{\sigma_0}(\num{0}, \num{n})
\land \dots \land \Obj S_{\sigma_k}(\num{k}, \num{n}) \land
\lforall[x][(\num{k} < x \lif \Obj S_\TMblank(x, \num{n}))]
\]
其中 $q$ 是 $M$ 在时刻~$n$ 的状态，$M$ 在时刻~$n$ 正在扫描方格~$m$，对于 $0 \le i \le k$，方格~$i$ 在时刻~$n$ 包含符号~$\sigma_i$，并且 $k$ 是时刻~$0$ 时纸带上最右边的非空白方格，或者 $n$ 步后读写头访问过的最右边的方格，取两者中较大的那个。
\end{defn}

\begin{lem}\ollabel{lem:halt-config-implies-halt}
如果 $M$ 在输入~$w$ 上运行 $n$ 步后处于一个停机配置，那么 $!C(M, w, n) \Entails !E(M, w)$。
\end{lem}

\begin{proof}
假设 $M$ 在输入~$w$ 上于 $n$ 步后停机。
则存在某个状态~$q$、方格~$m$ 和符号~$\sigma$，使得：
\begin{enumerate}
\item $n$ 步后，$M$ 处于状态~$q$，正在扫描包含~$\sigma$ 的方格~$m$。
\item 转移函数 $\delta(q, \sigma)$ 无定义。
\end{enumerate}
$!C(M, w, n)$ 描述了这个配置，并且将包含子句 $\Obj Q_{q}(\num{m}, \num{n})$ 和 $\Obj
S_{\sigma}(\num{m}, \num{n})$。这两个子句一起蕴涵 $!E(M, w)$：
\[
\lexists[x][\lexists[y][(\bigvee_{\tuple{q, \sigma} \in
      X}(\Obj Q_q(x, y) \land \Obj S_\sigma(x, y)))]]
\]
因为 $\Obj Q_{q'}(\num{m}, \num{n}) \land S_{\sigma'}(\num{m}, \num{n}) \Entails \bigvee_{\tuple{q, \sigma} \in X} (\Obj Q_q(\num{m}, \num{n}) \land \Obj S_{\sigma}(\num{m}, \num{n}))$，这是由于 $\tuple{q',\sigma'} \in X$。
\end{proof}

\begin{explain}
因此，如果 $M$ 在输入 $w$ 上停机，那么存在某个~$n$ 使得 $!C(M,
w, n) \Entails !E(M,w)$。我们现在来证明，对于任意时刻~$n$，都有 $!T(M, w) \Entails !C(M, w, n)$。
\end{explain}

\begin{lem}
\ollabel{lem:config}
对每个 $n$，如果 $M$ 在 $n$ 步后尚未停机，那么 $!T(M, w)
\Entails !C(M, w, n)$。
\end{lem}

\begin{proof}
归纳基始：若 $n = 0$，则 $!C(M, w, 0)$ 的合取项同时也是 $!T(M, w)$ 的合取项，因此被后者所蕴涵。

归纳假设：若 $M$ 在第 $n$ 步之前尚未停机，则 $!T(M,w) \Entails !C(M, w, n)$。我们需要证明（除非 $!C(M, w, n)$ 描述了一个停机配置），$!T(M, w) \Entails
!C(M, w, n+1)$。

假设 $n > 0$，且在 $n$ 步之后，以 $w$ 为输入启动的 $M$ 处于状态~$q$，正在扫描方格~$m$。由于 $M$ 在~$n$ 步后并未停机，其程序中必有一条指令属于以下三种形式之一：

\begin{enumerate}
\item \ollabel{right} $\delta(q, \sigma) = \tuple{q', \sigma', \TMright}$

\item \ollabel{left} $\delta(q, \sigma) = \tuple{q', \sigma', \TMleft}$

\item \ollabel{stay} $\delta(q, \sigma) = \tuple{q', \sigma', \TMstay}$
\end{enumerate}

我们将依次考虑这三种情况。

\begin{enumerate}
\item 假设有一条形如~\olref{right} 的指令。
  根据 \olref[rep]{defn:tm-descr}\olref[rep]{rep-right}，这意味着
\begin{align*} 
& \lforall[x][\lforall[y][((\Obj Q_{q}(x,
  y) \land \Obj S_\sigma(x, y)) \lif {}]]\\
& \qquad (\Obj Q_{q'}(x',
  y') \land \Obj S_{\sigma'}(x, y') \land !A(x, y)))  
 \intertext{是 $!T(M,w)$ 的一个合取项。这蕴涵了以下语句（全称实例化，以 $\num{m}$ 代~$x$，以 $\num{n}$ 代~$y$）：}
& (\Obj Q_{q}(\num{m}, \num{n}) \land \Obj S_{\sigma}(\num{m},
\num{n})) \lif {}\\
& \qquad (\Obj Q_{q'}(\num{m}', \num{n}') \land
\Obj S_{\sigma'}(\num{m}, \num{n}') \land !A(\num{m}, \num{n})).
 \intertext{由归纳假设，$!T(M, w) \Entails !C(M, w, n)$，即}
& \Obj Q_q(\num{m}, \num{n}) \land \Obj S_{\sigma_0}(\num{0}, \num{n})
\land \dots \land \Obj S_{\sigma_k}(\num{k}, \num{n}) \land \\
& \qquad\lforall[x][(\num{k} < x \lif \Obj S_\TMblank(x, \num{n}))]\\
 \intertext{由于在 $n$ 步后，纸带第 $m$ 个方格包含~$\sigma$，对应的合取项是~$\Obj S_\sigma(\num{m}, \num{n})$，因此这蕴涵：}
& \Obj Q_{q}(\num{m}, \num{n}) \land \Obj S_{\sigma}(\num{m},
   \num{n})
 \intertext{现在我们得到}
& \Obj Q_{q'}(\num{m}', \num{n}') \land \Obj S_{\sigma'}(\num{m},
  \num{n}') \land {}\\
& \qquad\Obj S_{\sigma_0}(\num{0}, \num{n}') \land \dots \land
  \Obj S_{\sigma_k}(\num{k}, \num{n}') \land {}\\
& \qquad \lforall[x][(\num{k} < x \lif \Obj S_\TMblank(x, \num{n}'))]
\end{align*}
具体如下：第一行直接由前一个条件句的后件经肯定前件得出。中间一行（不含 $S_{\sigma_m}(\num{m},\num{n}')$）的每个合取项，均由 $!C(M, w, n)$ 中的对应合取项连同 $!A(\num{m},
\num{n})$ 推出。

若 $m < k$，则 $!T(M,w) \Proves \num{m} < \num {k}$
（\olref[rep]{prop:mlessk}）且由~$<$ 的传递性，我们有
$\lforall[x][(\num{k} < x \lif \num{m} < x)]$。若 $m = k$，则
$\lforall[x][(\num{k} < x \lif \num{m} < x)]$ 仅凭逻辑即可得出。
于是最后一行可由 $!C(M, w, n)$ 的对应合取项、$\lforall[x][(\num{k} < x \lif \num{m} < x)]$ 以及 $!A(\num{m},
\num{n})$ 得出。若 $m<k$，这已经是 $!C(M, w, n+1)$。

现在假设 $m=k$。在这种情况下，在 $n+1$ 步后，读写头也访问了方格~$k+1$，该方格现已成为访问过的最右边的方格。
故而 $!C(M, w, n+1)$ 有一个新的合取项 $\Obj
S_\TMblank(\num{k}',\num{n}')$，且最后一个合取项是
$\lforall[x][(\num{k}' < x \lif \Obj S_\TMblank(x, \num{n}'))]$。我们需要验证这两个语句也被蕴涵。

我们已有 $\lforall[x][(\num{k} < x \lif \Obj S_\TMblank(x,
  \num{n}'))]$。特别地，由此可得 $\num{k} < \num{k}' \lif
\Obj S_\TMblank(\num{k}', \num{n}')$。从公理 $\lforall[x][x <
  x']$ 可得 $\num{k} < \num{k}'$。经肯定前件，得出 $\Obj
S_\TMblank(\num{k}',\num{n}')$。

另外，由于 $!T(M,w) \Proves \num{k} < \num{k}'$，关于~$<$ 的传递性公理给出 $\lforall[x][(\num{k}' < x \lif \Obj
  S_\TMblank(x, \num{n}'))]$。（此验证留作练习。）

\item 假设有一条形如~\olref{left} 的指令。
  那么，由 \olref[rep]{defn:tm-descr}\olref[rep]{rep-left}，
\begin{align*} 
& \lforall[x][\lforall[y][((\Obj Q_{q}(x', y) \land \Obj
    S_{\sigma}(x', y)) \lif {}]]\\
& \qquad (\Obj Q_{q'}(x, y') \land \Obj
  S_{\sigma'}(x', y') \land !A(x, y))) \land {}\\
& \lforall[y][((\Obj Q_{q_i}(\num{0}, y) \land \Obj S_{\sigma}(\num{0},
    y)) \lif {}]\\
& \qquad (\Obj Q_{q_j}(\num{0}, y') \land \Obj S_{\sigma'}(\num{0}, y')
  \land !A(\num{0}, y)))
 \intertext{是 $!T(M,w)$ 的一个合取项。若 $m>0$，则令 $l = m - 1$（即 $m = l+1$）。上述语句的第一个合取项蕴涵以下内容：}
& (\Obj Q_{q}(\num{l}', \num{n}) \land \Obj S_{\sigma}(\num{l}', \num{n}))
\lif {} \\
& \qquad (\Obj Q_{q'}(\num{l}, \num{n}') \land \Obj S_{\sigma'}(\num{l}',
\num{n}') \land !A(\num{l}, \num{n}))
 \intertext{否则，令 $l = m = 0$，并考虑第二个合取项所蕴涵的以下语句：}
& ((\Obj Q_{q_i}(\num{0}, \num{n}) \land \Obj S_{\sigma}(\num{0}, \num{n})) \lif {}\\
& \qquad   (\Obj Q_{q_j}(\num{0}, \num{n}') \land \Obj S_{\sigma'}(\num{0}, \num{n}') \land
!A(\num{0}, \num{n}))) 
 \intertext{任一语句均可如前推出}
&  \Obj Q_{q'}(\num{l}, \num{n}') \land \Obj S_{\sigma'}(\num{m},
  \num{n}') \land {}\\
&\qquad  \Obj S_{\sigma_0}(\num{0}, \num{n}')\land \dots \land
  \Obj S_{\sigma_k}(\num{k}, \num{n}')  \land {} \\
& \qquad  \lforall[x][(\num{k} < x
  \lif \Obj S_\TMblank(x, \num{n}'))]
\end{align*}。
（注意，第一种情形下 $\num{l}' \ident \num{l+1}
\ident \num{m}$，第二种情形下 $\num{l} \ident \num{0}$。）而这正是 $!C(M, w, n+1)$。

\item 情况 \olref{stay} 留作练习。
\end{enumerate}
我们已证明，对任意~$n$，$!T(M, w) \Entails !C(M, w, n)$。
\end{proof}

\begin{prob}
完成 \olref[tur][und][ver]{lem:config} 证明中的情况~\olref[tur][und][ver]{stay}。
\end{prob}

\begin{prob}
给出从
$\Obj S_{\sigma_i}(\num{i}, \num{n})$ 和 $!A(m, n)$ 到 $\Obj S_{\sigma_i}(\num{i}, \num{n}')$ 的一个推导（假设 $i \neq
m$，即 $i < m$ 或 $m < i$）。
\end{prob}

\begin{prob}
给出从 $\lforall[x][(\num{k} < x \lif \Obj
  S_\TMblank(x, \num{n}'))]$、$\lforall[x][x < x']$ 以及
$\lforall[x][\lforall[y][\lforall[z][ ((x < y \land y < z) \lif x <
      z)]]]$ 到 $\lforall[x][(\num{k}' < x \lif \Obj
  S_\TMblank(x, \num{n}'))]$ 的一个推导。
\end{prob}

\begin{lem}
\ollabel{lem:valid-if-halt}
若 $M$ 在输入~$w$ 上停机，则 $!T(M, w) \lif
!E(M, w)$ 是有效的。
\end{lem}

\begin{proof}
由 \olref{lem:config} 可知，对任意时刻~$n$，关于 $M$ 在时刻~$n$ 的配置描述~$!C(M, w, n)$ 被~$!T(M, w)$ 所蕴涵。
假设 $M$ 在 $k$ 步后停机。此时，存在某个 $m \in \Nat$，使得它正在扫描方格~$m$。
那么 $!C(M, w, k)$ 描述了~$M$ 的一个停机配置，即它同时包含 $\Obj Q_q(\num{m}, \num{k})$ 和 $\Obj
S_\sigma(\num{m}, \num{k})$ 作为其合取项，且 $\delta(q,\sigma)$ 无定义。
因此，由 \olref{lem:halt-config-implies-halt}，$!C(M, w, k) \Entails !E(M,
w)$。但由于 $!T(M, w) \Entails !C(M, w, k)$，我们有 $!T(M, w)
\Entails !E(M, w)$，故而 $!T(M, w) \lif !E(M, w)$ 是有效的。
\end{proof}

\begin{explain}
要完成我们断言的验证，我们还必须确立逆方向的蕴涵关系：若 $!T(M, w) \lif !E(M, w)$ 是有效的，则
$M$ 以输入~$w$ 启动时确实会停机。
\end{explain}

\begin{lem}
\ollabel{lem:halt-if-valid}
若 $\Entails !T(M, w) \lif !E(M, w)$，则 $M$ 在输入~$w$ 上停机。
\end{lem}

\begin{proof}
考虑 $\Lang L_M$-结构~$\Struct M$，其论域为~$\Nat$，它将 $\Obj 0$ 解释为~$0$，将 $\prime$~解释为后继函数，将 $<$~解释为小于关系，并将谓词 $\Obj Q_q$ 和~$\Obj S_\sigma$ 解释如下：
\begin{align*}
  \Assign{\Obj Q_q}{M} & =
\Setabs{\tuple{m, n}}{\begin{array}{ll}\text{以~$w$ 为输入启动，在~$n$ 步之后，}\\ \text{$M$ 处于状态 $q$ 并扫描方格~$m$}\end{array}} \\
\Assign{\Obj S_\sigma}{M} & = \Setabs{\tuple{m, n}}{\begin{array}{ll}
\text{以~$w$ 为输入启动，在 $n$ 步之后，}\\ \text{$M$ 的方格~$m$ 包含符号~$\sigma$}\end{array}}
\end{align*}
换句话说，我们构造结构~$\Struct{M}$，使得它逐步描述了 $M$ 以输入~$w$ 启动时的实际运行过程。
显然，$\Sat{M}{!T(M, w)}$。若 $\Entails !T(M, w) \lif !E(M, w)$，
则亦有 $\Sat{M}{!E(M, w)}$，即
\[
\Sat{M}{\lexists[x][\lexists[y][(\bigvee_{\tuple{q, \sigma} \in
      X}(\Obj Q_q(x, y) \land \Obj S_\sigma(x, y)))]]}.
\]
由于 $\Domain{M} = \Nat$，故必存在 $m, n \in \Nat$ 使得
$\Sat{M}{\Obj Q_q(\num{m}, \num{n}) \land \Obj S_\sigma(\num{m},
\num{n})}$ 对于某个使 $\delta(q, \sigma)$ 无定义的 $q$ 和~$\sigma$ 成立。
根据~$\Struct M$ 的定义，这意味着 $M$ 以输入~$w$ 启动 $n$ 步之后处于状态~$q$ 且正在读取符号~$\sigma$，而转移函数无定义，即 $M$~已经停机。
\end{proof}

\end{document}